\begin{problem}{}{}

    Sea ${\{ u_n \}}_{n \in \mathbb{N}_0}$ la sucesión definida por recurrencia como sigue:
    $u_0 = 2$, $u_1 = 4$ y $u_n = 4 u_{n-1} - 3 u_{n-2}$ con $n \in \mathbb{N}, n\ge 2$.

    \begin{enumerate}[(a)]
        \item Calcule $u_2$ y $u_3$.
        \item Probar que $u_n = 3^n + 1$, para todo $n \in \mathbb{N}_0$.
    \end{enumerate}

\end{problem}
\vspace{1ex}

Hagamos primero las evaluaciones puntuales y luego probemos la fórmula general:

\begin{enumerate}[(a)]
    \item Tenemos:

          \begin{tabular}{|l l r}
             & $u_2$ & \\
            $=$ & $4 u_1 - 3 u_0$ & (definición recursiva) \\
            $=$ & $4 \cdot 4 - 3 \cdot 2$ & (casos base) \\
            $=$ & $16 - 6$ & (aritmética) \\
            $=$ & $10$ & (aritmética) \\
          \end{tabular}

          Siguiendo con $u_3$:

          \begin{tabular}{|l l r}
             & $u_3$ & \\
            $=$ & $4 u_2 - 3 u_1$ & (definición recursiva) \\
            $=$ & $4 \cdot 10 - 3 \cdot 4$ (caso base y cálculo anterior) \\
            $=$ & $40 - 12$ & (aritmética) \\
            $=$ & $28$ & (aritmética) \\
          \end{tabular}

          Concluyendo así que $u_2 = 10 \text{ y } u_3 = 28$.

    \item Veamos ahora que $u_n = 3^n + 1, n \in \mathbb{N}_0$ es la fórmula general.
          Para esto, podemos usar el principio de inducción completa como sigue:

          \begin{itemize}
            \item \textbf{Casos Base:} Veamos que $u_0 = 2 = 1 + 1 = 3^0 + 1$ al mismo tiempo que $u_1 = 4 = 3 + 1 = 3^1 + 1$.
                  De esta forma vemos que ambos casos base se cumplen y concluimos esta parte.

            \item \textbf{Caso inductivo:} Supongamos que $u_h = 3^h + 1$ para todo $h$ tal que $0 \le h \le n$ para $n \in \mathbb{N}_{\ge 2}$.
                  Veamos ahora que la igualdad vale para $n+1$. Como sigue:

                  \begin{tabular}{|l l r}
                     & $u_{n+1}$ & \\
                    $=$ & $4 u_n - 3 u_{n-1}$ & (definición recursiva) \\
                    $=$ & $4 (3^n + 1) - 3 (3^{n-1} + 1)$ & (hipótesis inductiva) \\
                    $=$ & $4 \cdot 3^n + 4 - 3 \cdot 3^{n-1} - 3$ & (distributividad) \\
                    $=$ & $4 \cdot 3^n - 3^n + 1$ & (prop.\ de potencias y aritmética) \\
                    $=$ & $3 \cdot 3^n + 1$ & (álgebra) \\
                    $=$ & $3^{n+1} + 1$ & (prop.\ de potencias) \\
                  \end{tabular}

                  Y esta última igualdad es la que buscabamos probar. Con esto, finalizamos el caso inductivo.
          \end{itemize}

          Con esto, el principio de inducción completa nos permite concluir que la igualdad vale para todo entero no negativo.
\end{enumerate}
